\subsubsection*{The character and the two lattices on the retained family}
\label{subsec:marked-localization}

The adjoint comparison supplies the good row at the specified lift. Its contributions to the numerator are local: at each twisted place they use a particular Hecke character and the two integral lattices attached to it.

Fix a member of the cut family at scale \(x\), a lift \(j\), and
a twisted coefficient place \(w\mid\ell\). Let \(i\) be one of
the Hecke localizations at \((j,w)\), with coefficient integers
\(O\), and let \(\pi\) be a good eigencharacter there. In the
following argument the index \(i\) specifies the twist, its side
and its auxiliary level \(Q_i\). Write \(E\) for the auxiliary CM field
and \(v\) for a place of \(E\) above \(\ell\).

Let \(r\) be the rank of the complementary representation and \(t\) the
number of characters in the theta string. For the \(\ell\)-adic cyclotomic
character \(\chi_\ell\) of \(G_E=\operatorname{Gal}(\overline E/E)\), set
\begin{equation}\label{eq:signature-cyclotomic-minimum}
 \nu_\ell=\min_{g\in G_E}v_\ell\!\left(\chi_\ell(g)^{tr}-1\right),
 \qquad v_\ell(\ell)=1.
\end{equation}

\begin{proposition}[The signature and the complement at a marked point]
\label{prop:marked-signature}
Fix the norm choice on the retained family before choosing the
integral realizations at \(i\). The complement of the theta string
at \(\pi\) identifies, with its integral lattice, with one of the
good signatures tested on the prescribed side in
Proposition~\ref{prop:one-retained-character}. For the length \(s_\pi\) of
the annihilating quotient and the inertia readings \(\zeta_v\)
of the chosen twist, the simultaneous estimates are
\[
 s_\pi<\infty,\qquad s_\pi\le\nu_\ell+\ell^{-2},\qquad
 t\sum_{v\in Q_i} v_\ell(\zeta_v-1)\le\ell^{-2}.
\]
The one choice at \(w\) applies to all these presentations.
\end{proposition}

\begin{proof}
\proofstep{Read the normalized characteristic polynomial on the chosen side.}
Proposition~\ref{prop:cut-multiplicity} gives an occurring good
lower row \(\sigma\) for \(\pi\) at lift \(j\). In the coefficient model
take a nonzero common eigenvector \(x_\sigma\) in its good
isotypic space. The descent of an eigenvector here uses only
linear algebra. Take the basis of the small-field form, also a
basis after extension to the row field, and choose a nonzero
coordinate \(u\) of \(x_\sigma\). Over the subfield containing
the eigenvalues and the small-field matrices, choose a linear
functional on the row field with value \(1\) at \(u\).

Apply it
to each coordinate and then apply the coefficient embedding of
the chosen side. Write \(\lambda\) for the chosen functional and apply it coordinatewise. For every small-field Hecke matrix \(T\) and eigenvalue \(a_T\) in the field of linearity,
\[
 T\lambda(x_\sigma)=\lambda(Tx_\sigma)=a_T\lambda(x_\sigma),\qquad \lambda(u)=1.
\]
Applying the coefficient embedding preserves these identities. For the level action, undo the row character on each
group element. Its inverse value is in the eigenvalue field,
so the map also intertwines the sum giving the good projector.

The resulting vector lies in the good level space on that side.
At a split unramified prime \(P\) outside the finite omitted set,
its \(h\)-th Hecke eigenvalue is the embedded row eigenvalue
times \(c(P)^h\), with \(c(P)\ne0\) the normalizing character
factor in the level action. Decompose the good level space at \(j\) on this side into its recorded simultaneous eigenspaces. A nonzero component yields a signature \(s\) with the same normalized eigenvalues for every \(1\le h\le r\).

For corresponding unnormalized eigenvalues \(A_h,B_h\) and coefficient embedding \(\iota\), the cancellation reads
\[
 c(P)^h\iota(A_h)=c(P)^h\iota(B_h)\quad\Longrightarrow\quad A_h=B_h.
\]
For the normalized split roots \(\alpha_1,\ldots,\alpha_r\)
these values are the coefficients in
\[
 \prod_{i=1}^r (T-\alpha_i)
     =\sum_{h=0}^r(-1)^h e_h(\alpha_1,\ldots,\alpha_r) T^{r-h}.
\]
Thus equality of the embedded eigenvalues after dividing by
\(c(P)^h\) gives equality of polynomials, with all multiplicities. Undoing the common reciprocal and
scalar normalization identifies the Frobenius polynomials.
The signature is one of the good ones tested in
Section~\ref{subsec:retained-twist}: the localization at \(w\)
uses a nonzero twist index, and its good type and lower
cuspidal system are precisely those just selected.
\label{cor:marked-localization-1}

\proofstep{Choose integral realizations of all retained characters.}
Take the full semisimple integral Galois realization \(R_x\)
for each generic eigencharacter \(x\) of this localization,
as in Section~\ref{subsec:retained-galois}. At a theta
character the characteristic-zero representation has its
labelled string of characters and a continuous integral
complement \(\rho_x\). Denote the labelled characters by \(\chi_h\). At a good split Frobenius element \(g\), the comparison removes the string by the factorization
\[
 \det(1-XR_x(g))=\det(1-X\rho_x(g))\prod_h(1-X\chi_h(g)).
\]
The attached cuspidal summands realize the complementary factor. Density and semisimplicity identify the rational complement.

To obtain a stable lattice, start with any full lattice
\(L_0\) in the complement. Compactness gives a single integer
\(a\ge 0\) such that, for a coefficient uniformizer \(\pi\),
\[
 L_0\ \subseteq\ L=\sum_{g\in G_E}\rho_x(g)L_0
       \ \subseteq\ \pi^{-a}L_0,\qquad
 \rho_x(h)L=L\quad(h\in G_E).
\]
The upper lattice is Noetherian, so the middle module is a
full lattice. This construction only concerns the chosen complement.
The coefficient choice of Section~\ref{subsec:retained-galois}
splits the constituents and the roots of the marked corner.
At a retained non-theta character only \(R_x\) enters the matrix algebra. The unused complement may be assigned
\[
 \rho_x=\mathbf1^{\oplus r}\qquad(x\ \text{outside the theta locus}).
\]

This constructs one family at all the retained characters.
\label{cor:marked-localization-2}

\proofstep{Identify the two character evaluation kernels.}
For the comparison with \(s\), write \(O_0\) for the signature's
base coefficient ring. The character \(\psi_0\) on the parent
at \(w\) has its values in a common coefficient field. Its
two evaluations give the Galois character at \(i\) and the
character in the twist field, with the recorded side action.
The localization construction extends the embeddings to a
common field. On \(O_0\) it uses the prescribed map into
\(O\) and the inclusion into the twist integers. On the
parent field the two maps compare the same embedding at \(w\).
Concretely, two embeddings inducing \(w\) differ on the
completed algebraic coefficients by a local field embedding.
Extend over the finite coefficient fields into an algebraic
closure and transport the integral embedding by this
conjugacy.

For the local conjugacy \(\sigma\) and the absolute valuation \(v\), one has
\[
 v(\sigma a)=v(a)\ge0\qquad(a\in O_0).
\]
Thus the transported base coefficients are integral. The field used
for the localized row includes these base readings by
construction. Let \(\operatorname{ev}_i\) and \(\operatorname{ev}_s\) be the
two character evaluations on \(O_0[G_E]\), and let \(\iota_i,\iota_s\)
be their embeddings into the common field. For
\(a=\sum_g a_g[g]\) one has
\[
 \iota_i(\operatorname{ev}_i a)
     =\sum_g\iota_i(a_g)\iota_i(\operatorname{ev}_i[g])
     =\sum_g\iota_s(a_g)\iota_s(\operatorname{ev}_s[g])
     =\iota_s(\operatorname{ev}_s a).
\]
Injectivity of the common-field embeddings gives, for \(a\in O_0[G_E]\),
\[
 \operatorname{ev}_i(a)=0\iff\iota_i\operatorname{ev}_i(a)=0\iff\iota_s\operatorname{ev}_s(a)=0\iff\operatorname{ev}_s(a)=0.
\]
Applying this equality to \([g]-[1]\) also identifies
triviality on each decomposition group. The side inversion in taking the Galois
character of the ray class character is undone when the
signature is read back on that side.
\label{cor:marked-localization-3}

\proofstep{Transfer the length and inertia estimates on normalized lattices.}
The comparison of representations becomes a length estimate only after the character, the signature lattice and the chosen side of the place have been identified together.

Apply Lemma~\ref{lem:retained-signature-specialization}
to this \(\rho_\pi\) and this \(s\). The split polynomial identity
identifies the two rational rank-\(r\) representations with the
prescribed twists. After integral scaling the intertwiner \(A\) has nonzero reduction. Absolute residual irreducibility gives
\[
 \ker\bar A=0\quad\Longrightarrow\quad\det A\in O^\times,
\]
since the two lattices have the same rank. The equality of evaluation kernels
identifies the generated orders. Extension from their common
normalization preserves normalized length and the injection
into the signature Selmer group. The ramification set on the
signature side is exactly \(Q_i\). At each such place a marked inertia generator gives a generator
\(\zeta\) of the same finite character image. A replacement by
\(\zeta^a\) with \(a\) prime to its order preserves
\(v_\ell(\zeta-1)\).

Choose positive representatives \(a,b\) with \(ab\equiv1\) modulo
the order of \(\zeta\). The two divisibilities follow from
\[
 \zeta^a-1=(\zeta-1)\sum_{j=0}^{a-1}\zeta^j,\qquad
 \zeta-1=(\zeta^a-1)\sum_{j=0}^{b-1}\zeta^{aj}.
\]
Thus the two elements generate the same ideal. Thus its simultaneous character estimate gives the asserted
bounds on the actual complement.
The character used for this application was fixed with the
norm choice before \(R_x,\rho_x\) were selected. At the cut
the maps on the retained points restrict the original
readings, including \(Q_i\) and the side. This proves the
annihilation estimate for all of these integral presentations
with one choice at \(w\). \label{cor:marked-localization-4}
\end{proof}

At a labelled embedding of the host at \(v\), the archimedean
weights give the following candidate lists. Write \(c=(t-1)/2\), and use the real place \(\tau\)
under this embedding. In the Harish--Chandra parameter at
\(\tau\) put \(s=e_{j-1}\), write \(v_a\) for its \(k-1\)
decreasing offsets and \(b_h\), \(0\le h<m\), for the
increasing list interlaced with the \(e\)'s. On the positive
half the three sets for \(R_x\chi_1^{-1}\) are
\[
 C_+=\{c+s+v_a\}_a,\qquad D_+=\{c+s-b_h\}_h,\qquad
 I=\{0,\ldots,t-1\}.
\]
On the other half reflect the lists by \(a\mapsto t-1-a\),
writing \(C_-,D_-\) for the first two images.

\begin{lemma}[Signed jumps on the row and column]
\label{lem:marked-signed-jumps}
Writing \(\operatorname{HT}_d\) for the multiset of jumps on side \(d\) at a retained eigencharacter, one has
\[
 \operatorname{HT}_d(R_x\chi_1^{-1})=C_d\sqcup D_d\sqcup I,
\]
with all multiplicities one. At \(\pi\) the complement \(W=\rho_\pi\chi_1^{-1}\)
has jumps \(C_d\cup D_d\) and avoids \(0,-1\). At paired split
places the positive column of \(W\) and negative row of
\(W^\vee(1)\) each have \(m-j\) negative jumps per embedding.
\end{lemma}

\begin{proof}
\proofstep{Normalize the level-zero jump by the first character.}
Let \(p\) be a Harish--Chandra parameter
entry, \(d=\pm1\) the sign of the embedding and \(c_\tau\)
the common archimedean shift. The common shift cancels on twisting by \(\chi_1^{-1}\). Subtracting its jump from the level-zero jump gives
\[
 (dp+c_\tau)-(-ds-c+c_\tau)=d(p+s)+c.
\]

For the three pieces of the shifted parameter \(p+s\) the
entries are \(s+v_a\), \(s-b_h\), and the consecutive integers
from \(-c\) to \(c\). This proves the indicated reflection and
multiplicity formula. The cut gives the level-zero jump
identity on every retained eigencharacter, including at a
type prime. Away from the type primes its unramified
Weil--Deligne parameter and zero monodromy give
crystallinity; at a type prime use the potentially crystalline
comparison with the retained tame label.
\label{cor:marked-localization-5}

\proofstep{Separate the consecutive string from the signed complement.}
The inequalities on the selected weights give
\(0\le v_a\le S_c\), \(|s-b_h|\le\eta_c\), and
\(s>\eta_c+t+S_c\). The interlacing margins give the opposite signs on either side of \(j\):
\[
 s-b_h\ge(t+3)/2\ (h<j),\qquad s-b_h\le-(t+3)/2\ (h\ge j).
\]
These follow from the interlacing
margins on either side of \(e_{j-1}\) and monotonicity of
the \(b_h\). The three sets are disjoint and each jump is
simple. In particular removing \(I\) at \(\pi\) leaves
exactly \(C_d\cup D_d\) for \(W=\rho_\pi\chi_1^{-1}\);
this avoids \(0,-1\). The Tate dual \(W^\vee(1)\) at
the \emph{same} embedding sends each remaining \(a\) to
\(-1-a\).

For a complementary shifted entry \(z=p+s\), the positive
column jump and the negative row's dual jump are
\[
     a_+=z+c,\qquad a_{\rm row}=-1-(-z+c)=z-c-1 .
\]
For the \(b_h\)-entries with \(h<j\) one has \(z\ge c+2\);
both jumps are then positive. For \(h\ge j\) one has
\(z\le-c-2\), and both are negative. The outer entries
\(z=s+v_a\) exceed \(c+1\). Thus each side has \(m-j\)
negative jumps per embedding.
Thus each contributes at most the degree times \(m\) to
the free finite cohomology bound, even at the last lift.
\end{proof}

Above the prime cutoff the stronger separation required in
paragraph~\ref{subsec:outside-windows} follows directly from
the parameter inequality
\( |b_h-e_{j-1}|\ge 2b/\ell_*+t/2+3\). On either side
this absolute value is \(|a-c|\) for \(a\in D_d\).
On the complementary block \(C_d\) the place windows instead give
\[
 d(a-c)=s+v_a,\qquad W_0\le s\le W_0+\eta_c+1.
\]

Here \(W_0=b\) on the part with larger weights, \(\omega\) on the part with smaller weights and \(b_v+\lfloor b'/x\rfloor\) at a designated top.
The floor accounts for the extra one at the top. Thus the
exceptional integers tested there use the actual jump sets of
the marked map. These weight inequalities belong to the same
parameter record as \(\ell_*\).
\label{cor:marked-localization-6}

\subsubsection*{The two integral maps and the local cases}

Lemma~\ref{lem:marked-signed-jumps} determines the ordered jumps on the two marked maps. Proposition~\ref{prop:marked-signature} supplies the length on the good complement. The prime range determines which map to use.
Fix the exceptional primes computed from the host and the
three weight windows, and put the selected character and its
two auxiliary twist parameters in the family before choosing
the matrices \(R_x,\rho_x\). Choose the integral period functors
for the completions of this host with their comparison
isomorphisms, as in Sections~\ref{subsec:cut-wach-models}
and~\ref{subsec:split-local-length}.  In the rational matrix product
over the characters put the fibre at \(\pi\) in the direct
sum basis for the string and its complement.

If this changes
basis by \(B\) over the fraction field, it carries every sum
\(\sum_g a_g R(g)\), with coefficients in the localized
Hecke order, to the sum with the same coefficients in the new
generators. Indeed
\[
 B\Big(\sum_g a_gR(g)\Big)B^{-1}
      =\sum_g a_g\,BR(g)B^{-1}.
\]
Applying the identity with \(B^{-1}\) gives the reverse inclusion of
the integral spans even if \(B\) is defined only over the fraction field. This transports the marked
idempotents, the ordered corner roots and the element
perturbing the cyclotomic corner. On the rank-\(r\) block
the integral span maps onto \(M_r(O)\): its image modulo the
maximal ideal is the whole matrix algebra by absolute residual
irreducibility, and the cokernel vanishes by Nakayama.
Section~\ref{subsec:split-selmer} carries this out with the
split Frobenius polynomials of the chosen localization.

Write \(T\) for the stable lattice of \(W\) in this direct sum.

\begin{lemma}[Lattices for the local row and column comparisons]
\label{lem:marked-row-column-lattices}
In the integral presentation the evaluation at the rank-one corner gives a surjection
\[
 e_1\mathcal A\twoheadrightarrow T^*=\operatorname{Hom}(T,O)(1).
\] The rational column
representations have constituents in sums of \(R_x\chi_1^{-1}\).
Their relations give finite classes on the rank-\(r\) target,
an unramified class on the first character line, and zero on
the other character lines.
\end{lemma}

\begin{proof}
\proofstep{Use the full matrix image for the character row.}
Write the rank-one projection as \(v\lambda\) with
\(v\in T\), \(\lambda\in\operatorname{Hom}_O(T,O)\), and
\(\lambda(v)=1\). For any \(\lambda'\in\operatorname{Hom}_O(T,O)\),
choose the matrix \(M=v\lambda'\) in the full matrix image.
The row of this product, evaluated on any \(x\in T\), is determined by
\[
 ((v\lambda)M)x=v\lambda\!\left(v\lambda'(x)\right)
                =v\lambda(v)\lambda'(x)=v\lambda'(x).
\]
Thus the image contains every integral row.
Put \(U_g=\chi_1(g)^{-1}R(g)\). The action on the row functional, including the Tate factor, is
\[
 g\cdot\lambda'=\epsilon(g)\lambda'\circ U_g^{-1}\quad\text{on }T^*.
\]
The source embeds rationally into a sum
of the dual-twisted family by recording all its matrix rows.
\proofstep{Use the direct-block projection for the column lifts.}
The row argument used the entire matrix image. On the conjugate column the finite classes are read from the projection onto the original complementary block.

For the column lifts one uses the left action on matrix
columns and projection to the direct block. Its rational
constituents belong to sums of \(R_x\chi_1^{-1}\).
The column relations give finite classes in the rank-\(r\)
block, an unramified class on the first character line and
zero on the other lines, as shown in
Section~\ref{subsec:split-selmer}.

For the character-row map
the extension with trivial quotient is integral: on
\(Y=\sum_g\mathcal A(U_g-1)E\) use the cocycle \((U_g-1)E\).
The map
\[
 (z,a)\mapsto(z+aE,a)
\]
embeds its rationalization
in the matrix columns and a trivial line. The latter already
occurs in the twisted string. These constructions justify the
passage of the jump bounds to the lifting maps by exactness
on crystalline subquotients. The modulus multiplying an
integral finite class is computed on the actual surjection.
For a finite extension \(O\subset O'\) of coefficient DVRs and a finite-length module \(M\),
\[
 \frac{\operatorname{length}_{O'}(M\otimes_O O')}{e(O')}=\frac{e(O'/O)\operatorname{length}_O(M)}{e(O'/O)e(O)}=\frac{\operatorname{length}_O(M)}{e(O)}.
\]
The normalized colength
returns to this presentation. \label{cor:marked-localization-7}
\end{proof}

Apply Lemma~\ref{lem:marked-row-column-lattices} to the
following cases on this same marked family.

\paragraph{Unramified split and inert places.}
For split places choose the side whose embedding sign makes
the direct column positive, and take the conjugates of these
places for its rows. If \(\mathcal R_\ell\) and \(\mathcal C_\ell\) denote the row and column places, at an inert prime these sets are
\[
 \mathcal R_\ell=\{v:v\mid\ell\},\qquad \mathcal C_\ell=\varnothing.
\]
The unramified Wach estimates apply to these surjections, with the order of Frobenius on the embeddings of \(E\) and the signed jump arrays above.
Sections~\ref{subsec:ordinary-split-length} and~\ref{subsec:inert-rows}
give the resulting split and inert allowances. Above the
cutoff, Paragraph~\ref{subsec:outside-windows} applies to
this same family with the exceptional set already specified.
Proposition~\ref{prop:offwindow-local-vanishing} kills the additional
local allowance outside these exceptions.

\paragraph{The split conditions at two and three.}
At two, and at three when \(3\nmid Q_S\), partition the places into columns and conjugate rows using the archimedean sign. One choice covers every good \(\pi\) at the localization and preserves the matrices, twist, support and character-row bound.
The column and row lifting exponents are proved for these crystalline surjections over
\(\mathbb Q_\ell(\mu_{\operatorname{lcm}(\ell,N)})\) and
returned by corestriction in Section~\ref{subsec:split-local-length}.
The preliminary choice of zero for an unused exponent supplies
no such lifting result. The character-row estimate and its dual use the respective heights
\[
 h=\lfloor2b\rfloor,\qquad h^\vee=h+1,\qquad n_0=4nR.
\]
The extra unit allows the Tate-dual shift.
The depth is at least two at \(2\). At \(3\mid Q_S\) use the
ramified odd-prime estimate of
Section~\ref{subsec:special-prime-lengths}.
This includes the torsion of finite integral cohomology when
its rational rank is zero. \label{cor:marked-localization-8}

\paragraph{The common denominator of the two type projectors.}
At a type prime apply the pair of finite series projectors
for the labels of \(\pi\), multiplied by their common group
denominator. If \(e_\pi\) denotes their product in the rational
finite group algebra and \(a_\pi=h_Ge_\pi\) its integral
multiple by the common group denominator \(h_G\), then
\[
       a_\pi|_{\operatorname{im} e_\pi}=h_G\,\mathrm{id},
        \qquad a_\pi|_{\ker e_\pi}=0 .
\]
Both projectors commute with the Hecke action since the
level actions and the split Hecke correspondences commute.
The remaining support has the two inertia labels of \(\pi\)
at its theta points. Such a point is good by the cut.

The
norm choice identifies the inertia characters on this
support with the labels of the same finite series, including
inversion on the conjugate side. Hence the common tame
lattice comparison of Section~\ref{subsec:special-prime-lengths}
applies on the support of this operator. Use
\[
 \eta=\lfloor\eta_c\rfloor,\ s_T=\lfloor3b\rfloor+2,\ S_T=0
\]
with the opposite signs at the paired places and \(C_T=40\).
The auxiliary prime bound there gives the strict size inequality
once \(b\ge C_T\). The pushout estimate then bounds the
finite-index defect for the rank-\(r\) target and each of
its character targets. Multiplying the resulting element
of the congruence order by that very projector clears its
complement. This gives the type-prime length on each graded
piece. \label{cor:marked-localization-9}

\paragraph{A threshold shared by the local length estimates.}
All these estimates may be used at a single late point.
The unramified row comparison first fixes its absolute
constant \(C_2\ge1\). The type-prime, small-prime, parameter
and exceptional-prime comparisons supply finitely many
further constants. Write \(A_1,\ldots,A_s\) for the further upper bounds after insertion of \(C_2\). Choose
\[
 C\ge\max\{C_2,A_1,\ldots,A_s\}.
\]

For
that bound take the maximum of the thresholds for the
integral setup, the size of the exceptional set, the
special primes, the large-prime vanishing, the ordinary
split and inert estimates, the parameter inequalities,
and the bounds at two and three. Each has already been
obtained for all point indices above its own threshold.
Only now take the point and evaluate its exceptional set
and integral models. An empty good locus has zero projected
lattice. At each other localization the rank-\(r\) term has
the signature bound above and the character terms have
the strict and class-number bounds of
Section~\ref{subsec:division-length}. This gives the local
table used in the weighted return on precisely these cuts.
\label{cor:marked-localization-10}
